\section*{Chapter \ref{ch:g6:food-and-energy} --- Food and Energy}

\begin{solution}{exo:g6:food-and-energy:1}
About $1.4$ tonnes per hectare in 1961 and about $4.2$ in 2024 (the table gives
$1.35$ and $4.21$).
\end{solution}

\begin{solution}{exo:g6:food-and-energy:2}
$36 \div 8 = 4.5$ tonnes per hectare.
\end{solution}

\begin{solution}{exo:g6:food-and-energy:3}
United States ($6\,454$), China ($2\,947$), Germany ($2\,874$), India ($777$), Chad
($171$ kilograms of oil equivalent per person).
\end{solution}

\begin{solution}{exo:g6:food-and-energy:4}
$600 \div 2\,400 \times 100\,\% = 25\,\%$.
\end{solution}

\begin{solution}{exo:g6:food-and-energy:5}
(a) $4.21 \div 1.35 \approx 3.1$ times. (b) $(4.21 - 1.35) \div 1.35 = 2.86 \div 1.35
\approx 2.12$: a rise of about $212\,\%$.
\end{solution}

\begin{solution}{exo:g6:food-and-energy:6}
(a) In the early 1970s (the table gives 1973). (b) About 30 million tonnes then,
and about $1.5$ million tonnes in 2023.
\end{solution}

\begin{solution}{exo:g6:food-and-energy:7}
(a) A: $400 \div 1\,000 = 40\,\%$; B: $800 \div 4\,000 = 20\,\%$. (b) B spends more on
food (800 against 400 dalens); A has the larger \omterm{def:g6:food-and-energy:food-share}{food share}. (c) The richer family
spends more on food, but a smaller share of its income.
\end{solution}

\begin{solution}{exo:g6:food-and-energy:8}
(a) The lines of the United States, the United Kingdom and Germany: these countries
were much richer in 2023 than in 1970, and emitted far less sulphur dioxide. (b) The
text gives carbon dioxide for one year only (2023): richer countries then emitted
more per person, but that does not say what happens when one country gets richer.
One would need each country's carbon dioxide year after year, beside its income.
\end{solution}

\begin{solution}{exo:g6:food-and-energy:9}
(a) $6\,676 \div 2\,874 \approx 2.3$ times. (b) Almost the same distance from the
left (same income); Saudi Arabia's dot is higher, by a little more than a third of
a gridline step.
\end{solution}

\begin{solution}{exo:g6:food-and-energy:10}
(a) Now: harvest $100 \times 3.0 = 300$ tonnes, need $1\,000 \times 0.3 = 300$ tonnes.
After 10 years: $100 \times 3.5 = 350$ against $1\,200 \times 0.3 = 360$. After 20:
$400$ against $420$. After 30: $450$ against $480$. (b) From the first decade on:
10 tonnes short after ten years, 30 tonnes short after thirty. (c) The need grows by
$200 \times 0.3 = 60$ tonnes every ten years, so the harvest of 100 hectares must grow
by 60 tonnes: the \omterm{def:g6:food-and-energy:crop-yield}{yield} must rise by $60 \div 100 = 0.6$ tonne per hectare every ten
years.
\end{solution}

\begin{solution}{exo:g6:food-and-energy:11}
(a) $(37.85 - 10.82) \div 37.85 = 27.03 \div 37.85 \approx 0.714$: a fall of $71.4\,\%$.
(b) A country's emissions of an \omterm{def:g6:food-and-energy:air-pollutant}{air pollutant} can fall a great deal while it keeps
growing richer.
\end{solution}

\begin{solution}{exo:g6:food-and-energy:12}
(a) For example: laws limiting the sulphur in fuels; filters fitted to power
stations; coal replaced by gas, nuclear or wind; factories that moved abroad.
(b) No: the chart shows only the totals. (c) For example, for filters: the
number of power stations with filters, year by year, compared with the fall; for
factories moving abroad: the country's imports of steel and cement, and the
emissions of the countries that make them.
\end{solution}

\begin{solution}{pb:g6:food-and-energy:1}
\textbf{1.} $300 \div 40 = 7.5$ tonnes per hectare.
\textbf{2.} $40 \times 6.0 = 240$ tonnes; $(300 - 240) \div 240 = 0.25$, a rise of
$25\,\%$.
\textbf{3.} $300 \times 1\,500 = 450\,000$ loaves, enough for $450\,000 \div 200 = 2\,250$
people for a year.
\textbf{4.} Poorest $3\,989 \div 14\,314 \approx 27.9\,\%$; second $4\,516 \div 19\,334
\approx 23.4\,\%$; third $4\,823 \div 24\,228 \approx 19.9\,\%$; richest $5\,251 \div
35\,720 \approx 14.7\,\%$.
\textbf{5.} $5\,251 - 3\,989 = 1\,262$ dalens more.
\textbf{6.} The richer the quarter, the more dalens it spends on food, but the
smaller the share of its spending that goes to food.
\textbf{7.} New total $14\,314 \times 1.1 = 15\,745.4$ dalens; $3\,989 \div 15\,745.4
\approx 25.3\,\%$.
\textbf{8.} $120\,000 \times 0.02 = 2\,400$ tonnes.
\textbf{9.} $60\,000 \times 0.02 = 1\,200$ tonnes.
\textbf{10.} $1\,200 \times (1 - 0.90) = 120$ tonnes.
\textbf{11.} $(2\,400 - 120) \div 2\,400 = 0.95$: a fall of $95\,\%$. The town paid
for the gas, which may cost more than coal, and for the filter and its upkeep;
electricity may have become dearer.
\textbf{12.} For example: ``The richer households spend more on food but a smaller
share of their budget. A richer town can afford cleaner fuels and filters, so it can
have less smoke even while it uses as much energy.''
\end{solution}
